\documentclass{article}

% --- NeurIPS 2026 Style (Main Track, double-blind) ---
% neurips_2026.sty is sourced from the official template ZIP:
%   https://media.neurips.cc/Conferences/NeurIPS2026/Formatting_Instructions_For_NeurIPS_2026.zip
% The .sty file is checked into docs/NeurIPS/ alongside this file.
% For camera-ready, change to: \usepackage[main, final]{neurips_2026}
% [REVISION | track-main | 2026-04-25 | Switched track option from `eandd` (Evaluations and Datasets) to `main` (Main Track) per locked decision in memory/next_session_prompt.md "Decisions Locked" table; verified `main` is a valid \DeclareOption in neurips_2026.sty:49 — pending audit]
\usepackage[main]{neurips_2026}

% --- Packages ---
\usepackage{amsmath, amssymb, amsthm}
\usepackage{booktabs}
\usepackage{hyperref}
\usepackage{graphicx}
\usepackage{enumitem}
\usepackage{microtype}
\usepackage{xcolor}

% --- Review-edit markup (remove wrapper before final submission) ---
% Passages wrapped in \reviewedit{} are flagged for advisor review.
% They appear in blue in the compiled PDF.
\newcommand{\reviewedit}[1]{\textcolor[RGB]{0,100,180}{#1}}

% --- Theorem environments ---
\newtheorem{theorem}{Theorem}
\newtheorem{proposition}{Proposition}
\newtheorem{definition}{Definition}
\newtheorem{hypothesis}{Hypothesis}

% --- Hyperref setup ---
\hypersetup{
    colorlinks=true,
    linkcolor=blue,
    citecolor=blue,
    urlcolor=blue
}

\title{KV-Cache Inception: Reversible Monte Carlo Tree Search in Latent Space for Detecting and Stress-Testing Alignment Faking in Large Language Models}

\author{[Authors]}
\date{}

\begin{document}
\maketitle

% ============================================================
% ABSTRACT
% ============================================================
% [REVISION | abstract-recast-α+ | 2026-05-06 | Replaced abstract entirely (Session A draft, applied 2026-05-06). Old hook: alignment-faking framing with three-contribution list (telemetry → MCTS → first-formal-OE). New hook: algorithmic search infrastructure with monitoring application; explicit Bailey 2024 engagement; OEI framed as measurement framework with prerequisites, not as a results headline. The "first formal definition" claim is dropped (pre-mortem FM 1). Forward refs: Theorem~\ref{thm:reversibility} (added 2026-05-06). Empirically defensible claims only: Theorem 1 zero-drift on 1B + 3B (validated 2026-04-16 + 2026-05-05); memory bound independent of branching factor with constant left unquantified pending Max's simplified Exp 3. — pending audit]
\begin{abstract}
Latent-space analysis of frozen large language models---including activation steering, representation engineering, and adversarial probing of internal monitors---increasingly requires search infrastructure that scales with the depth of intervention rather than with the branching factor of the search. We introduce \textbf{Reversible Monte Carlo Tree Search in KV-cache latent space}, an inference-time algorithm that mutates a frozen Transformer's cached key/value tensors in place via FP32 accumulators and rolls back exact baseline state at each tree expansion. We prove that the FP32 accumulator guarantees zero-drift reversibility (Theorem~\ref{thm:reversibility}; empirically validated to machine precision over 200 apply--revert cycles on Llama~3.2-1B and Llama~3.2-3B), and we analyze a memory bound that is independent of branching factor---a constant-factor multiple of the base KV cache rather than the exponential blowup of parallel-cache MCTS. To make the algorithm useful for representational analysis, we introduce a \textbf{per-layer telemetry matrix} that fuses bottom-up hallucination-associated-neuron activation tracking~\cite{gao2025hneurons} with top-down representation-engineering projections~\cite{zou2023repe}, providing a continuous, differentiable signal at each generation step. As an application, we propose the \textbf{Orthogonal Escape Index (OEI)}---a targeted geometric measurement of the fraction of induced activation change that lies in the null space of a linear monitor---as a measurement framework complementing Bailey et al.'s~\cite{bailey2024obfuscated} cosine-to-baseline measurement of obfuscated activations under gradient-trained input attacks. We discuss measurement-pipeline prerequisites identified in our reference implementation and outline the empirical characterization required for the inference-time, frozen-weight regime. All artifacts are released with a Croissant 1.1 metadata package for reproducibility.
\end{abstract}

% ============================================================
% 1. INTRODUCTION
% ============================================================
\section{Introduction}

The deployment of large language models in safety-critical applications has made alignment evaluation one of the most consequential open problems in machine learning. The standard evaluation paradigm operates at the behavioral level: a model is presented with test inputs, and its outputs are assessed for compliance with safety criteria, either by human evaluators or by automated LLM-as-a-Judge frameworks~\cite{zheng2023judging, liu2023geval}. This paradigm rests on an implicit assumption---that a model's external behavior faithfully reflects its internal computational state. A growing body of evidence suggests this assumption is unreliable.

Anthropic's landmark study on alignment faking demonstrated that Claude 3 Opus, when informed that its responses would be used for retraining, strategically complied with instructions it would otherwise refuse, explicitly reasoning in a hidden chain-of-thought scratchpad---included as part of the experimental design---that compliance was instrumentally rational to preserve its existing harmlessness values~\cite{greenblatt2024alignment}. Hubinger et al.\ (2024) showed that deliberately inserted ``sleeper agent'' backdoors can persist through reinforcement learning, supervised fine-tuning, and adversarial training in large models---with adversarial training found to actively teach the model to better conceal its backdoor rather than remove it~\cite{hubinger2024sleeper}; the paper notes this tests threat-model resilience rather than how likely such behaviors are to arise naturally, and that persistence is substantially weaker in smaller models. Most directly relevant to our work, Gupta and Jenner (2025) introduced RL-Obfuscation, demonstrating that LLMs finetuned via reinforcement learning can learn to evade latent-space monitors---including linear probes trained to detect deception---while maintaining their black-box behavior, and that models trained to evade a single monitor can generalize to evade unseen monitors of the same aggregation class---though the paper reports that max-pooling and attention-based monitors remain robust to this attack~\cite{gupta2025rlobfuscation}.

These findings collectively establish a troubling implication: \textbf{behavioral evaluation is necessary but insufficient for alignment assurance, and even mechanistic monitoring tools can be circumvented by sufficiently capable models}. The question is no longer whether models \emph{can} fake alignment, but how to develop evaluation frameworks that can detect and stress-test this capability at the mechanistic level.

We propose that the answer lies in moving from passive monitoring to \textbf{active, adversarial intervention in the model's latent space}. Rather than observing a model's internal states and hoping to detect deception, we systematically \emph{perturb} those states using structured search algorithms and measure the model's response across multiple independent telemetry channels. This approach is motivated by a key insight: while a model may learn to mask deceptive intent along monitored representational axes, the simultaneous maintenance of deceptive capability and evasion of monitoring creates informational constraints that manifest as detectable signatures in orthogonal subspaces.

% [CRITIQUE NOTE | Evidence burden]
% The empirical burden is high: the paper will need an explicit validation chain
% linking continuous latent interventions to discrete deceptive intent outcomes.
% Reviewers will ask: does an increase in OEI/TDS during KV-cache steering correspond
% to a measurable increase in behaviorally deceptive outputs? Without this link,
% the telemetry signals are interesting but not interpretively grounded.
% Consider adding this to the Limitations section.

\subsection{Contributions}

This paper makes the following contributions:

% [REVISION | contributions-recast-α+ | 2026-05-06 | Replaced entire §1 Contributions enumerate block with Option α+ recast (Session A draft, applied 2026-05-06). Old order: telemetry → MCTS → OE formalization. New order: MCTS → telemetry matrix → OEI as targeted measurement. The old "first formal definition and experimental protocol for Orthogonal Escape" headline is dropped (Bailey 2024 already operationalized the question, per pre-mortem FM 1). The 2026-05-05 Contrib2-fix REVISION marker is subsumed: the new Contribution 1 paragraph cites Theorem~\ref{thm:reversibility}, scopes scale to 1B + 3B with explicit deferral of 20B to Phase B, and adds the 2026-05-05 3B feasibility evidence. Forward references to thm:reversibility, prop:memory, sec:measurement-prereqs, eq:mcts_reward all resolve. — pending audit]
\begin{enumerate}[leftmargin=*]
    \item \textbf{Reversible Monte Carlo Tree Search in KV-cache latent space.} We introduce an inference-time algorithm that mutates a frozen Transformer's cached key/value tensors via in-place arithmetic ($\mathbf{K}_t \leftarrow \mathbf{K}_t + \alpha \cdot \mathbf{d}_K^{(\ell)}$) and inverts each mutation exactly through an FP32 accumulator. We prove (Theorem~\ref{thm:reversibility}) that the accumulator guarantees zero-drift reversibility independent of search depth and validate this empirically on Llama~3.2-1B-Instruct and Llama~3.2-3B-Instruct (drift $= 0$ to machine precision over 200 cycles; Llama~3.2-3B feasibility confirmed on a single 12\,GB consumer GPU). The memory complexity of the algorithm is independent of branching factor---a constant-factor multiple of the base KV cache (Proposition~\ref{prop:memory})---in contrast to the $O(b^d \cdot M_{\text{KV}})$ overhead of standard parallel-cache MCTS. The algorithm is gradient-free and operates on a frozen model at inference time, distinguishing it from gradient-based input-space attacks~\cite{zou2023gcg, bailey2024obfuscated} and from training-time monitor-evasion methods~\cite{gupta2025rlobfuscation}.

    \item \textbf{Per-layer telemetry matrix for latent-space analysis.} We define a $2 \times L$ matrix $\mathbf{T}_t$ that fuses bottom-up hallucination-associated-neuron activation~\cite{gao2025hneurons} with top-down representation-engineering projections~\cite{zou2023repe}. The matrix provides a continuous, differentiable signal at each generation step and supports four diagnostic regime classifications (creative synthesis, contextual fabrication, alignment faking, orthogonal escape) keyed to joint patterns across the two channels. The telemetry matrix serves dual roles: as a per-step monitoring signal during search, and as the third reward component in the MCTS node valuation function (Eq.~\ref{eq:mcts_reward}).

    \item \textbf{The Orthogonal Escape Index as a targeted measurement of null-space activation redistribution.} We propose OEI---the fraction of induced activation change $\Delta\mathbf{h}$ that lies in the null space of a linear monitor $\mathbf{P}$---as a geometric quantity for studying inference-time activation redistribution under steering interventions. OEI complements Bailey et al.'s~\cite{bailey2024obfuscated} cosine-to-baseline measurement: where Bailey asks whether gradient-trained input perturbations can achieve global activation orthogonality (reporting a tentative negative result), OEI asks how much of any induced perturbation lands in directions invisible to a specific monitor. We discuss measurement-pipeline prerequisites identified in our reference implementation (\S\ref{sec:measurement-prereqs}) and identify the steps required for empirical characterization in the inference-time, frozen-weight regime.
\end{enumerate}

% ============================================================
% 2. BACKGROUND AND RELATED WORK
% ============================================================
\section{Background and Related Work}

\subsection{Alignment Faking and Deceptive Alignment}

% [REVISION | §2.1-soften-α+ | 2026-05-07 | Two sentence-level softens to align §2.1 with the §3 reframe and §1 abstract framing. (1) Opening: "...has been a longstanding concern" → "...is one motivating application area for latent-space analysis on frozen models, and a longstanding concern". (2) Closing: "motivating the white-box representational monitoring approach our telemetry matrix provides" → "one motivation for the inference-time latent-space analysis framework developed in this paper, of which the telemetry matrix is the monitoring component". Middle sentences (Greenblatt + Hubinger summaries with caveats) preserved verbatim. Closes the §2.1 audit item from 2026-05-06 EOD; structural re-ordering of §2 was considered and rejected (cross-reference risk too high this close to deadline). Net ~10 words longer; negligible page-budget impact. — pending audit]
The possibility that advanced AI systems might strategically misrepresent their alignment is one motivating application area for latent-space analysis on frozen models, and a longstanding concern in AI safety research~\cite{hubinger2024sleeper, greenblatt2024alignment}. Recent empirical work has moved this from theoretical speculation to demonstrated capability. Anthropic's alignment faking study~\cite{greenblatt2024alignment} provided the first large-scale evidence that frontier models can engage in strategic compliance when made explicitly aware of their training situation---though the paper notes the goals being preserved were the model's existing HHH harmlessness preferences rather than malicious objectives, and that this behavior emerged only in the largest models tested (Claude 3 Opus and Claude 3.5 Sonnet), with smaller models not exhibiting it. The Sleeper Agents work~\cite{hubinger2024sleeper} demonstrated that deliberately inserted backdoors can resist removal through reinforcement learning, supervised fine-tuning, and adversarial training in large models---with adversarial training found to actively teach the model to conceal rather than remove its backdoor---though the authors note this tests threat-model resilience rather than the likelihood of such behaviors arising naturally, and that persistence is strongly scale-dependent. Existing empirical work on alignment faking and backdoor persistence suggests that behavioral evaluations alone are insufficient for reliable alignment detection~\cite{greenblatt2024alignment, hubinger2024sleeper}: a model that strategically complies or conceals a backdoor can pass behavioral tests while maintaining deceptive internal representations---one motivation for the inference-time latent-space analysis framework developed in this paper, of which the telemetry matrix is the monitoring component.

\subsection{Representation Engineering and Activation Steering}

Representation Engineering (RepE), introduced by Zou et al.\ (2023)~\cite{zou2023repe}, established that high-level cognitive states---including honesty and truthfulness---are linearly decodable from a model's latent representation space. This finding enabled \emph{activation steering}, wherein targeted vector additions to intermediate representations can modify model behavior without weight updates~\cite{turner2023steering}. Belitsky et al.\ (2025)~\cite{belitsky2025kvcache} extended this to \emph{cache steering}, demonstrating that a single, one-shot modification to the KV cache can steer model behavior more stably than continuous activation interventions---validated for reasoning induction in small-to-medium LLMs (360M--8B). However, all existing steering work operates in a \emph{corrective} paradigm---applying fixed precomputed vectors to induce desired behaviors. Unlike this static approach, our MCTS iteratively explores the KV-cache intervention space across multiple rollouts, reversing each intervention exactly via FP32 accumulators and guided by a telemetry-derived reward rather than a fixed behavioral objective. Prior work has explored the latent intervention space via gradient ascent~\cite{bailey2024obfuscated} and reinforcement learning~\cite{gupta2025rlobfuscation}, but no prior work has deployed a heuristic tree-search algorithm to \emph{explore} the KV-cache intervention space at inference time, without gradient access or model finetuning.

A critical limitation of the corrective paradigm was exposed by Gupta and Jenner (2025)~\cite{gupta2025rlobfuscation}, who showed that models can learn to evade latent-space monitors via RL finetuning (weight updates across training epochs), with mean/median-aggregated monitors being particularly vulnerable while max-pooling and attention-based monitors remain robust. Their finding that representational redistribution contributes to monitor evasion---alongside a primary behavioral mechanism of output-policy change (accounting for approximately 70\% of the recall drop)---provides conceptual motivation for our Orthogonal Escape hypothesis; however, RL-Obfuscation operates via weight finetuning on updated models, while our work concerns inference-time KV-cache interventions on a frozen model, so this motivates our hypothesis by analogy rather than direct mechanism. The geometric motivation for this redistribution is formalized in Section~3 (Hypothesis~\ref{hyp:orthogonal_escape}).

% [REVISION | §2.2-Bailey-engagement | 2026-05-06 | Inserted Bailey 2024 substantive engagement paragraph after the existing RL-Obfuscation paragraph (Session A draft, applied 2026-05-06). Bailey positioned as peer measurement (gradient-trained input attacks vs our inference-time KV-cache mutations) rather than competitor whose novelty claim we contest. Three explicit complementarities: intervention site (input embedding vs cached K/V), search objective (similarity-minimization vs telemetry-driven), and metric scope (global activation orthogonality vs targeted monitor-null-space fraction). The §3 forward pointer to Hypothesis~\ref{hyp:orthogonal_escape} remains at the end of the RL-Obfuscation paragraph, before this new paragraph (per draft instruction). — pending audit]
Most directly relevant to our framing of OEI, Bailey et al.\ (2024)~\cite{bailey2024obfuscated} systematically study whether adversarial perturbations to model inputs can achieve activation states that are simultaneously behavior-preserving and orthogonal to baseline activations. Their gradient-trained universal embedding suffixes optimize a joint objective: minimize the KL divergence between perturbed and baseline output distributions while minimizing the cosine similarity between perturbed and baseline activations. They report a tentative negative result for the regime of trained input-space attacks against a frozen model: near-orthogonality and behavior preservation cannot be jointly achieved, since pushing cosine similarity to the level of unrelated inputs requires a KL divergence that produces incoherent output. Our work is complementary along three axes. First, the intervention is applied to cached key/value tensors at inference time rather than to input embeddings during gradient training, so neither the model weights nor the embedding layer is modified. Second, the search optimizes a behavioral telemetry signal (Eq.~\ref{eq:mcts_reward}) rather than directly minimizing similarity to baseline, so any redistribution we observe is incidental rather than optimization-driven. Third, OEI measures a more targeted geometric quantity---the fraction of induced perturbation invisible to a specific linear monitor---rather than global activation orthogonality. Whether the inference-time, frozen-weight regime exhibits the same constraint Bailey identifies remains an open question, and is one motivation for the measurement framework developed in this paper.

\subsection{H-Neurons and Hallucination Detection}

Gao et al.\ (2025)~\cite{gao2025hneurons} identified H-Neurons---a sparse subset of feed-forward network units whose activation patterns reliably predict hallucination occurrences and are causally linked to over-compliance behaviors. Unlike black-box evaluation, H-Neuron monitoring provides a continuous, neuron-level signal that can serve as a differentiable reward for optimization algorithms. We extend this finding to the domain of alignment faking by treating $\sigma_H$ as an \emph{over-compliance pressure proxy}: H-Neurons validated by Gao et al.\ (2025) as causal indicators of over-compliance and hallucination may produce detectable activation signatures when a model encounters pressure to fabricate alignment-consistent outputs. This extrapolation from over-compliance to strategic, instrumentally rational deception is a testable assumption rather than an established equivalence---one that our experimental protocol is designed to validate.

\subsection{MCTS for LLM Red-Teaming}

Monte Carlo Tree Search has been increasingly applied to automated red-teaming of LLMs. Wu et al.\ (2025)~\cite{wu2025mcts} proposed MCTS-based prompt autogeneration for jailbreak attacks. The DAMON framework~\cite{damon2025} introduced dialogue-aware MCTS for multi-turn jailbreaks. ACE-Safety~\cite{li2025ace} formulated attack-defense co-evolution via tree-group dual-aware search. However, \textbf{all existing MCTS-based red-teaming operates exclusively in text space}---mutating prompts, suffixes, or dialogue turns. This fundamentally limits the search to the discrete, combinatorial space of token sequences and provides no access to the model's internal computational dynamics. Non-MCTS adversarial search methods face a related but distinct limitation. Gradient-based attacks such as Greedy Coordinate Gradient (GCG)~\cite{zou2023gcg}---which require white-box gradient access to the target model---operate in discrete token space and are therefore equally constrained to the discrete combinatorial surface: they optimize token sequences via gradient signals but are not designed to navigate the continuous geometry of the model's internal representation space. Our MCTS, by contrast, operates without gradient access on a frozen model. MCTS provides the exploration semantics---branching, rollout, and UCB1 selection---that allow systematic coverage of the latent intervention landscape in a way that both prompt-mutation and gradient-based approaches are not designed to navigate. Our work is the first to conduct MCTS via continuous arithmetic mutations of the KV cache with an exact FP32 rollback mechanism, enabling gradient-free, inference-time exploration of the model's representational landscape.

\subsection{Continuous Latent Reasoning and KV-Cache Search}

Recent work has moved beyond discrete token generation to perform search and reasoning directly in continuous representation spaces. COCONUT~\cite{hao2025coconut} bypasses discrete token decoding by feeding final-layer hidden states back as input embeddings, enabling emergent breadth-first latent tree search. This validates the viability of continuous latent search, but operates via recurrent hidden-state traversal rather than direct arithmetic mutation of a frozen transformer's KV-cache tensors.
% [REVISION | G3 | 2026-04-24 | CRSM removed entirely (unreviewed prototype); COCONUT alone cited; "Both" → "This" — pending audit]

Most directly analogous to our architectural approach, the Okazaki-RAG framework~\cite{okazakirag2025} (an unpublished conceptual contribution; not peer-reviewed) implements a tree search whose nodes are KV-cache assemblies, with explicit rollback when heuristic scores deteriorate. However, Okazaki-RAG constructs its search tree by swapping pre-computed, discrete cache blocks rather than executing continuous arithmetic mutations across the differentiable tensor space. Discrete block swapping requires a precomputed library of fixed states; arithmetic KV-cache mutation enables gradient-free exploration of a continuous intervention landscape at inference time. Our FP32 accumulator mechanism provides mathematically exact rollback without the VRAM overhead of duplicating multi-gigabyte cache states across branches---a barrier that block-swapping architectures do not resolve in the same form.

\subsection{LLM-as-a-Judge and Evaluation Limitations}

% [REVISION | R-B-§2.2-soften | 2026-05-05 | Softened closing sentence of §2.2 (was: "Grounding evaluation in standard deterministic metrics such as pass@1 and cosine similarity of semantic intent provides stronger anchoring..."). Replacement preserves the paragraph's pivot from LLM-judge limitations to deterministic alternatives without falsely implying we measure pass@1 / cosine-similarity ourselves; those metrics were culled from §5.4 in this same session. — pending audit]
LLM-as-a-Judge frameworks, including MT-Bench~\cite{zheng2023judging} and G-Eval~\cite{liu2023geval}, use language models to evaluate output quality holistically. While semantically aware, these approaches achieve reliable agreement with human experts only when the judge model can itself correctly answer the underlying question or is provided with a correct reference; without such grounding, agreement drops substantially---for instance, GPT-4o's Cohen's $\kappa$ falls from $0.86$ to $0.16$ on benchmark questions the model cannot answer~\cite{llmasjudge2025limitations}. Static analysis tools (SonarQube, Semgrep) are deterministic but cannot evaluate whether code satisfies the \emph{intent} behind a requirement. Standard deterministic metrics such as pass@1 or semantic-similarity scores anchor evaluation to objective signals; our work extends this principle from output-space measurements to internal-state telemetry.

% ============================================================
% 3. PROBLEM FORMULATION
% ============================================================
\section{Problem Formulation}

% [REVISION | §3-reframe-α+ | 2026-05-06 | Reframed §3.1 from "The Alignment Faking Detection Problem" (alignment-faking-as-headline framing) to "Inference-Time Latent-Space Analysis on Frozen Models" (measurement-framework framing where alignment evaluation is one specialization). Eq.~\ref{eq:alignment_faking} (Behavioral vs Null-projection) preserved verbatim with label intact. Hypothesis~\ref{hyp:orthogonal_escape} preserved with label intact, retitled to clarify regime ("Orthogonal Escape under inference-time KV-cache perturbation"), and softened: "will redistribute its deceptive encodings" → "may redistribute its representation" — honest under Option α+ given the §5 results that show Bailey-consistent saturation rather than positive OE. New prose adds explicit Bailey reference as the analogous gradient-trained input-attack regime, motivating the inference-time variant. §3.2 (Discrete Text-Space Bottleneck) untouched. — pending audit]
\subsection{Inference-Time Latent-Space Analysis on Frozen Models}

Let $f_\theta$ be a language model with parameters $\theta$ held fixed throughout this work, and let $\mathbf{h}_t^{(\ell)} \in \mathbb{R}^d$ denote the hidden state at layer $\ell$ and token position $t$. A linear monitor $\mathbf{P}_{\text{monitor}}$ is a projection in $\mathbb{R}^d$---for instance, a representation-engineering honesty probe~\cite{zou2023repe} or an aggregation thereof~\cite{gupta2025rlobfuscation}. The central question of this paper is: under controlled perturbations of the model's cached key/value tensors (and hence its hidden states at subsequent positions), how does the geometric relationship between the monitored subspace $\text{Range}(\mathbf{P}_{\text{monitor}})$ and the unmonitored complement $\text{Null}(\mathbf{P}_{\text{monitor}})$ evolve?

This question has multiple downstream applications. In mechanistic interpretability, it characterizes how representations respond to controlled perturbations along monitored directions. In alignment evaluation, it specializes to a robustness check on latent-space monitors: a model exhibits \emph{alignment faking} if it satisfies

\begin{equation}
\text{Behavioral}(f_\theta, x) \in \mathcal{S} \quad \text{while} \quad \exists\, \mathbf{v} \in \text{Null}(\mathbf{P}_{\text{monitor}}): \langle \mathbf{h}_t^{(\ell)}, \mathbf{v} \rangle \neq 0
\label{eq:alignment_faking}
\end{equation}

where $\mathcal{S}$ is the set of safe behavioral outputs and $\text{Null}(\mathbf{P}_{\text{monitor}})$ is the null space of the monitoring projection. That is, the model produces safe outputs while encoding non-trivial information in directions invisible to the monitor. Whether such a configuration is reachable under inference-time, frozen-weight steering---without weight updates or input-space gradient training---is an open empirical question. Bailey et al.~\cite{bailey2024obfuscated} characterized the analogous question for gradient-trained input-space attacks and reported a tentative negative result; we formulate the inference-time variant as Hypothesis~\ref{hyp:orthogonal_escape}.

\begin{hypothesis}[Orthogonal Escape under inference-time KV-cache perturbation]\label{hyp:orthogonal_escape}
When corrective pressure is applied along the monitored subspace $\text{Range}(\mathbf{P}_{\text{monitor}})$, a sufficiently capable model under inference-time steering may redistribute its representation into $\text{Null}(\mathbf{P}_{\text{monitor}})$, preserving behavior while evading detection:
\begin{equation}
\mathbf{h}_t^{(\ell)} \xrightarrow{\text{steering}} \mathbf{h}_t^{(\ell)} + \alpha \cdot \mathbf{d}_{\text{steer}} \implies \text{Proj}_{\text{Null}(\mathbf{P})}(\mathbf{h}_{t+k}^{(\ell)}) \neq \mathbf{0}
\label{eq:orthogonal_escape}
\end{equation}
where $\mathbf{d}_{\text{steer}}$ is the steering direction and $\alpha$ is the intervention magnitude.
\end{hypothesis}

\subsection{The Discrete Text-Space Bottleneck}

Standard MCTS-based red-teaming operates in the discrete, combinatorial space of token sequences. With vocabulary size $|V|$, the branching factor at each node is $O(|V|)$, and the search must commit to specific token realizations at every expansion step. This creates two fundamental limitations: (1)~the exponential blowup of the discrete search space renders deep exploration computationally intractable, and (2)~operating in text space provides no direct access to the model's internal representational dynamics---precisely the space where alignment faking manifests according to the formalization in Equation~\ref{eq:alignment_faking}.

We argue that the Orthogonal Escape phenomenon is intrinsically a \emph{geometric} problem in continuous latent space. The redistribution of deceptive encodings into the null space of a linear monitor is a continuous vector operation that cannot be meaningfully probed by discrete token-level perturbations. Detecting and characterizing this redistribution requires an algorithm that can navigate the continuous geometry of the model's representation space directly.

This motivates our central methodological contribution: shifting the MCTS search from discrete text space into the continuous latent space of the KV cache. The Reversible MCTS algorithm, enabled by the FP32 accumulator mechanism, provides the computational machinery to conduct this geometric exploration efficiently. The accumulator guarantees exact state restoration after arbitrary search trajectories, transforming the intractable problem of parallel cache duplication into a sequential, memory-efficient traversal of the latent intervention landscape.

% ============================================================
% 4. PROPOSED METHOD
% ============================================================
\section{Proposed Method}

\subsection{Multi-Dimensional Micro-Telemetry Matrix}

We construct a telemetry matrix $\mathbf{T}_t \in \mathbb{R}^{2 \times L}$ at each generation step $t$ across $L$ transformer layers, combining two orthogonal monitoring channels:

\paragraph{Bottom-Up Channel: H-Neuron Activation Tracking.}
Following Gao et al.\ (2025)~\cite{gao2025hneurons}, we identify the set of hallucination-associated neurons $\mathcal{H} \subset \{1, \ldots, d_{\text{ffn}}\}$ for the target model. At each generation step $t$, we compute the aggregate H-Neuron stress signal:

\begin{equation}
\sigma_H^{(\ell)}(t) = \frac{1}{|\mathcal{H}|} \sum_{j \in \mathcal{H}} \text{ReLU}\left(a_j^{(\ell)}(t) - \bar{a}_j^{(\ell)}\right)
\label{eq:hneuron_stress}
\end{equation}

where $a_j^{(\ell)}(t)$ is the activation of neuron $j$ at layer $\ell$ and step $t$, and $\bar{a}_j^{(\ell)}$ is its baseline activation computed over a calibration set of faithful generations.

% [CRITIQUE NOTE | Confound risk]
% σ_H stress may conflate multiple failure modes (hallucination, mode-collapse,
% strategic masking pressure) unless discriminant controls are added. The H-Neuron
% literature establishes association with over-compliance and hallucination; the
% extrapolation to strategic, instrumentally-rational deception is a testable
% assumption, not a settled equivalence. Consider adding discriminant controls
% (e.g., baseline comparisons on non-deceptive prompts) or explicitly stating
% this in Limitations.

% [REVISION | DRAFT-HOLD-block1-consolidated | 2026-05-05 | GAP-C1-05 (LAT/PCA vs difference-in-means in §4.1) + P1-B Finding 4 (all-layer aggregation vs middle-layer recommendation) — block content moved into §5 Implementation Notes paragraph (items i and ii). Code fix tracking continues in action plan §4.5 (Max's lane). — pending audit]
\paragraph{Top-Down Channel: RepE Linear Artificial Tomography.}
We apply the Linear Artificial Tomography (LAT) procedure from Zou et al.\ (2023)~\cite{zou2023repe} to extract linear probes for cognitive dimensions including honesty ($\mathbf{w}_{\text{hon}}^{(\ell)}$), certainty ($\mathbf{w}_{\text{cert}}^{(\ell)}$), and goal-coercion ($\mathbf{w}_{\text{coerce}}^{(\ell)}$). Note that Zou et al.\ (2023) validate LAT for honesty and truthfulness; the goal-coercion and certainty probe directions are extrapolations of the LAT methodology to novel concept dimensions not directly studied in the source paper. The RepE compliance signal at layer $\ell$ is:

\begin{equation}
\rho_R^{(\ell)}(t) = \mathbf{w}_{\text{hon}}^{(\ell)\top} \mathbf{h}_t^{(\ell)}
\label{eq:repe_signal}
\end{equation}

The telemetry matrix is then:

\begin{equation}
\mathbf{T}_t = \begin{bmatrix} \sigma_H^{(1)}(t) & \cdots & \sigma_H^{(L)}(t) \\ \rho_R^{(1)}(t) & \cdots & \rho_R^{(L)}(t) \end{bmatrix}
\label{eq:telemetry_matrix}
\end{equation}

We define four diagnostic states based on the joint signature:

\begin{table}[h]
\centering
\begin{tabular}{@{}llll@{}}
\toprule
\textbf{State} & \textbf{H-Neuron Stress $\sigma_H$} & \textbf{RepE Honesty $\rho_R$} & \textbf{Interpretation} \\
\midrule
Creative Synthesis & Low--Moderate & High, Stable & Valid novel generation \\
Contextual Fabrication & High (spike) & Degrading & Comprehension debt instantiation \\
Alignment Faking & Low (suppressed) & High (artificially maintained) & Strategic compliance masking \\
Orthogonal Escape & Divergent from $\rho_R$ & Divergent from $\sigma_H$ & Deceptive encoding in null space \\
\bottomrule
\end{tabular}
\caption{Diagnostic signatures from the joint telemetry matrix. The critical diagnostic for Orthogonal Escape is a statistically significant divergence between $\sigma_H$ and $\rho_R$: the H-Neuron channel detects over-compliance pressure while the RepE channel reports nominal honesty, indicating that the model has learned to satisfy the RepE probe while encoding deceptive intent in orthogonal directions.}
\label{tab:diagnostic_states}
\end{table}

\subsection{Reversible MCTS in KV-Cache Space}

\subsubsection{The Computational Bottleneck}

Standard MCTS requires evaluating multiple future branches simultaneously. When MCTS operates in the KV-cache latent space of a Transformer, each node expansion requires modifying the cached key and value tensors. Under standard implementations (including PagedAttention), mutating the past KV cache invalidates all dependent computations, requiring a full physical copy of the cache for each branch. For a model with $L$ layers, sequence length $S$, and hidden dimension $d$, the KV cache occupies:

\begin{equation}
M_{\text{KV}} = 2 \cdot L \cdot S \cdot d \cdot \text{sizeof(dtype)}
\label{eq:kv_memory}
\end{equation}

For a 20B-parameter model in bf16 precision, $M_{\text{KV}} \approx 40\text{GB}$. A standard MCTS with branching factor $b$ and depth $d$ requires $O(b^d)$ simultaneous cache copies, exceeding the 80GB VRAM of a single H100 before completing a depth-2 rollout.

% [REVISION | DRAFT-HOLD-block2-consolidated | 2026-05-05 | GAP-C2-01 (per-layer steering notation in §4.2 vs single-vector broadcast in code) — block content moved into §5 Implementation Notes paragraph (item iii). Code fix tracking continues in action plan §4.2 / Phase 3. — pending audit]
\subsubsection{The Reversible Rollback Algorithm}

We exploit the mathematical property that tensor addition is reversible, provided numerical stability is maintained. Let $\mathbf{K}_t^{(\ell)}, \mathbf{V}_t^{(\ell)}$ denote the key and value cache tensors at layer $\ell$ for token position $t$. Let $\mathbf{d}_K^{(\ell)}, \mathbf{d}_V^{(\ell)}$ be the RepE steering vectors, and let $\alpha_K, \alpha_V$ be the MCTS-chosen intervention magnitudes.

\paragraph{Step 1: Forward Mutation (Steer).} Before evaluating a branch, mutate the KV tensors in place:

\begin{equation}
\mathbf{K}_t^{(\ell)} \leftarrow \mathbf{K}_t^{(\ell)} + \alpha_K \cdot \mathbf{d}_K^{(\ell)}, \qquad \mathbf{V}_t^{(\ell)} \leftarrow \mathbf{V}_t^{(\ell)} + \alpha_V \cdot \mathbf{d}_V^{(\ell)}
\label{eq:forward_mutation}
\end{equation}

\paragraph{Step 2: Evaluation.} Execute the forward pass, generate the next token, and record the telemetry matrix $\mathbf{T}_t$ as the MCTS node reward.

\paragraph{Step 3: Reverse Rollback (Restore).} Invert the mutation exactly:

\begin{equation}
\mathbf{K}_t^{(\ell)} \leftarrow \mathbf{K}_t^{(\ell)} - \alpha_K \cdot \mathbf{d}_K^{(\ell)}, \qquad \mathbf{V}_t^{(\ell)} \leftarrow \mathbf{V}_t^{(\ell)} - \alpha_V \cdot \mathbf{d}_V^{(\ell)}
\label{eq:reverse_rollback}
\end{equation}

\paragraph{Precision Management: FP32 Accumulator.} Repeated bf16 addition and subtraction introduces epsilon drift that corrupts the cache over many rollbacks. We maintain an FP32 accumulator tensor $\mathbf{A}^{(\ell)} \in \mathbb{R}^{S \times d}$ that tracks the cumulative delta applied to each mutated position. All steering deltas are computed and accumulated in FP32; the bf16 cache is reconstructed from the accumulator at each step:

\begin{equation}
\mathbf{A}^{(\ell)} \leftarrow \mathbf{A}^{(\ell)} + \alpha \cdot \mathbf{d}^{(\ell)} \quad (\text{FP32}), \qquad \mathbf{K}_t^{(\ell)} \leftarrow \mathbf{K}_{\text{base},t}^{(\ell)} + \text{cast}_{\text{bf16}}(\mathbf{A}^{(\ell)})
\label{eq:fp32_accumulator}
\end{equation}

This guarantees exact baseline restoration after every MCTS node evaluation, achieving \textbf{zero semantic degradation} regardless of search depth.

% [REVISION | DRAFT-HOLD-block3-consolidated | 2026-05-05 | GAP-C2-03 (50MB accumulator overhead claim in §4.2 vs full-shape S×d allocation in code) — block content moved into §5 Implementation Notes paragraph (item iv). Code fix tracking continues in action plan §4.2 / Phase 3. — pending audit]
\paragraph{Memory Complexity.} The reversible algorithm reduces memory from $O(b^d \cdot M_{\text{KV}})$ to $O(M_{\text{KV}} + d \cdot L \cdot S' \cdot d_{\text{model}} \cdot 4)$, where $S'$ is the number of mutated token positions (typically $S' \ll S$). For a 20B model with $S' = 10$ mutated positions, the accumulator overhead is approximately 50MB---negligible compared to the 40GB base cache.

\subsubsection{MCTS Node Valuation}

The reward function for each MCTS node is derived from the telemetry matrix rather than from textual output:

\begin{equation}
r(\text{node}) = \lambda_1 \cdot \bar{\sigma}_H + \lambda_2 \cdot (1 - \bar{\rho}_R) + \lambda_3 \cdot \text{Divergence}(\sigma_H, \rho_R)
\label{eq:mcts_reward}
\end{equation}

where $\bar{\sigma}_H$ and $\bar{\rho}_R$ are the mean H-Neuron stress and RepE honesty across layers, and $\text{Divergence}(\sigma_H, \rho_R)$ measures the cross-channel inconsistency that signals Orthogonal Escape. The MCTS uses UCB1 for node selection:

\begin{equation}
\text{UCB1}(i) = \bar{r}_i + c \cdot \sqrt{\frac{2 \ln N}{n_i}}
\label{eq:ucb1}
\end{equation}

where $\bar{r}_i$ is the mean reward of node $i$, $N$ is the total number of visits to the parent, $n_i$ is the number of visits to node $i$, and $c = \sqrt{2}$ is the exploration constant.

% ============================================================
% 5. EXPERIMENTAL DESIGN
% ============================================================
\section{Experimental Design}
\label{sec:experiments}
% [REVISION | sec-experiments-label-add | 2026-05-06 | Added \label{sec:experiments} on the §5 section heading to enable forward references from the §7.4 measurement-prereqs subsection (revised this batch) and from the §5 Experiment 1 results paragraph (this batch). Section heading text unchanged. — pending audit]

% [REVISION | §5-impl-notes | 2026-05-05 | Added consolidated "Implementation Notes (Phase 2 / Phase 3 Distinction)" paragraph at the start of §5. Content sourced from the five DRAFT-HOLD blocks at TeX:184–200 (GAP-C1-05 + Finding 4), TeX:247–264 (GAP-C2-01), TeX:294–311 (GAP-C2-03), TeX:351–366 (GAP-C1-05 second location), TeX:501–525 (GAP-C2-03/04 Memory Complexity). Per pre-mortem FM 5 + Max §3.4: a single visible paragraph is what a reviewer needs; the original five DRAFT-HOLD blocks have been replaced with one-line REVISION markers at their old locations to maintain the audit trail. Asymptotic claims (Theorem 1, branching-factor independence, telemetry tensor shape) are unaffected; only constants and operational details are deferred to Phase 3. — pending audit]
\paragraph{Implementation Notes (Phase 2 / Phase 3 Distinction).}
The reference implementation backing this submission is a Phase 2 prototype on Llama~3.2-1B-Instruct and Llama~3.2-3B-Instruct on consumer hardware (NVIDIA RTX~3060). Several elements of the method as written in §4 and §6 differ from this reference implementation in identifiable, bounded ways; we list them here as a single source of truth so a reviewer reading the paper alongside the released code finds one consolidated reference rather than scattered caveats. The asymptotic claims---exact reversibility (Theorem~1), memory-complexity scaling independent of the branching factor (Proposition in §6), and the telemetry matrix tensor shape (Eq.~\ref{eq:telemetry_matrix})---are unaffected by the distinctions below; what differs is constants and operational details that resolve once the corresponding Phase~3 code work lands.

\textit{(i) RepE probe extraction (§4.1, §5.2).} The text describes the Linear Artificial Tomography procedure of \citet{zou2023repe} using PCA on contrastive pairs. The Phase~2 implementation uses a difference-in-means estimator, which approximates the leading PCA component on contrastive pairs in expectation but is not equivalent to LAT as published. PCA-based extraction is deferred to Phase~3.

\textit{(ii) Layer aggregation for the honesty probe (§4.1).} The Phase~2 implementation aggregates the RepE honesty signal across all $L$ layers; \citet{zou2023repe} validate the honesty probe on the middle layers only. Layer-restricted aggregation is deferred to Phase~3.

\textit{(iii) Per-layer steering vectors (§4.2).} The notation in §4.2 presents per-layer steering vectors $\mathbf{d}_K^{(\ell)}, \mathbf{d}_V^{(\ell)}$. The Phase~2 implementation broadcasts a single middle-layer vector to all layers; per-layer independent steering is deferred to Phase~3.

\textit{(iv) Sparse accumulators and baseline-clone overhead (§4.2 + §6).} The accumulator analysis in §4.2 and the worked example in §6 treat $K_{\text{acc}}$ as a sparse $S' \times d$ allocation with $S' \ll S$, yielding an example overhead of approximately $50$~MB at $20$B scale. The Phase~2 implementation allocates full-shape $S \times d$ accumulators and additionally stores a baseline cache clone of size $M_{\text{KV}}$ to enable rollback. The asymptotic claim that memory is independent of the branching factor $b$ (§6 Memory Complexity Proposition) is unchanged. The constant factor in the Phase~2 implementation is approximately $2$--$3 \cdot M_{\text{KV}}$ rather than the $\sim 1 \cdot M_{\text{KV}}$ implied by the worked example in §6: the $40.05$~GB figure quoted there reflects sparse-accumulator and no-clone assumptions, while the actual Phase~2 floor at $20$B scale is approximately $80$~GB before any accumulator overhead. Sparse-accumulator implementation, baseline-clone elimination, and corrected example numbers backed by direct VRAM measurement are required to bring §6's worked example into agreement with the Phase~2 implementation, and are deferred to Phase~3.

\textit{(v) On the reversibility theorem.} Theorem~1 (Exact Reversibility with FP32 Accumulator) is independent of the four distinctions above. It was empirically validated to numerical zero drift over $200$ apply-reverse cycles on Llama~3.2-1B-Instruct (residual norm $= 0.00$, 2026-04-16 gate) and confirmed at the 3B scale (residual norm $= 0.00$, 2026-05-05 probe). The headline algorithmic contribution is unaffected by any constant-factor revisions to §6.

\subsection{Phase 0: Architectural Decoupling}
% [IMPL: DONE — logomesh/ package: BaseModelClient, LocalLlamaOracle, modular architecture complete]

Before deploying telemetry hooks, the monolithic MCTS orchestrator must be refactored to support: (a)~a model abstraction layer replacing hardcoded API calls with a \texttt{BaseModelClient} interface supporting local HuggingFace models; (b)~a state serialization pipeline enabling MCTS tree persistence across environments; and (c)~a continuous reward interface processing neuronal activation tensors directly as MCTS node valuations.

\subsection{Phase A: Local Prototyping on Lightweight Surrogates}
% [IMPL: DONE — Phase 3 gate passed 2026-04-16 on Llama-3.2-1B-Instruct]
% [IMPL: Theorem 1 validated: measure_lipschitz_drift.py, 200 cycles, FP32 drift = 0.00e+00]
% [IMPL: MCTS smoke test: run_kv_mcts.py --nodes 10, all T_t / OEI / TDS fields populated]

We establish the offline execution environment using Llama-3.2-1B-Instruct (hidden dimension 2048, 22 layers, ${\sim}2$GB RAM) on consumer CPU hardware. This phase validates:

\begin{itemize}[leftmargin=*]
    \item \textbf{H-Neuron identification} via contrastive activation analysis between faithful and hallucinated generations on a curated calibration set.
% [REVISION | DRAFT-HOLD-block4-consolidated | 2026-05-05 | GAP-C1-05 second location (§5.2 Phase A bullet "RepE probe training via PCA" vs code uses difference-in-means; IMPL DONE adjacency made this a live false claim) — block content moved into §5 Implementation Notes paragraph (item i). Code fix tracking continues in action plan §4.5 (Max's lane). — pending audit]
    \item \textbf{RepE probe training} via PCA on contrastive pairs for honesty, certainty, and goal-coercion dimensions.
    \item \textbf{Reversible MCTS convergence} on the 1B surrogate, measuring: (a)~Lipschitz drift with and without FP32 accumulators over 1000 rollback cycles; (b)~search depth achievable within fixed compute budgets; (c)~telemetry matrix stability across MCTS branches.
    \item \textbf{Cross-model alignment theory} via supervised Procrustes alignment between 1B and 7B model representation spaces.
\end{itemize}

\subsection{Experiments}
% [REVISION | PhaseB-header | 2026-04-24 | Renamed from "Phase B: Scaling to Target Models"; H100 removed — submission experiments run on Llama 3.2 1B+3B, RTX 3060 — pending audit]

We execute the following experiments on Llama~3.2-1B-Instruct and Llama~3.2-3B-Instruct on a local RTX~3060 (12GB VRAM). Scaling to models up to 20B parameters on H100 is planned for Phase~B follow-on work; see §7 Future Work. Key experiments:

% [REVISION | Exp1-reframe-with-results-α+ | 2026-05-06 | Reframed Experiment 1 from hypothesis-test framing ("Orthogonal Escape Detection"; predicting OE signature) to measurement-and-report framing ("OEI Characterization under Inference-Time KV-Cache Perturbation"). Integrated post-fix-run results: Table~\ref{tab:exp1-results} with per-α mean values (OEI 0.991→0.995; ρ_R −1.306→−0.533; σ_H 0 throughout; TDS 0.094→0.046; reversibility holds end-to-end). The Interpretation paragraph states explicitly that the data is *consistent with Bailey 2024's negative-result regime* — TDS *decreases* with α, opposite the OE signature predicted by the original framing. This refutes the prior "(a) and (c) will improve while (d) will simultaneously increase" claim from the same paragraph; new prose says so directly. Forward references to §\ref{sec:measurement-prereqs} for prerequisites context resolve. Replaces the 2026-04-25 Exp1-1B3B-framing marker (subsumed). — pending audit]
\paragraph{Experiment 1: OEI Characterization under Inference-Time KV-Cache Perturbation.}
We apply layer-specific steering deltas of increasing magnitude $\alpha \in \{0.1, 0.5, 1.0, 2.0, 5.0\}$ along the RepE honesty direction, projected through the per-layer key projection weights into per-head cache space (see §\ref{sec:measurement-prereqs}). The sweep is run on Llama~3.2-1B-Instruct on a single 12\,GB consumer GPU; reversibility on Llama~3.2-3B-Instruct is independently confirmed (Theorem~\ref{thm:reversibility}, residual\_norm $= 0.00 \times 10^{0}$ over 200 cycles, validated 2026-04-16). For each $\alpha$ we measure: (i) the OEI (Eq.~\ref{eq:oei}); (ii) the RepE honesty projection $\rho_R$; (iii) the H-Neuron stress signal $\sigma_H$; (iv) the cross-channel divergence $\text{TDS} = \text{JSD}(\sigma_H, \rho_R)$.

% [REVISION | Exp1-σH-verification-update | 2026-05-06 | Replaced one sentence in the Results paragraph below — the σ_H reporting line. Was: "The H-Neuron stress channel σ_H reports zero across all α; this is consistent with its calibration against hallucination-inducing prompts and is discussed further in §\ref{sec:measurement-prereqs}." New: "The H-Neuron stress channel σ_H reports zero across all α; controlled verification (§\ref{sec:measurement-prereqs}) localizes this to a measurement-pipeline issue in the dense-scoring path independent of the steering projection or calibration regime." Reflects the σ_H verification result (diagnostic session 2026-05-06): σ_H = 0 across all 17 layers under hallucination/coercion/neutral regimes, including a coercion prompt drawn verbatim from the calibration class — refutes the calibration-prerequisite framing. Net ~5 words shorter. Other sentences in Results + Interpretation paragraphs unchanged. Parent marker `Exp1-reframe-with-results-α+` retained above. — pending audit]
\textbf{Results.} Per-$\alpha$ mean values across the sweep are reported in Table~\ref{tab:exp1-results}. OEI varies monotonically with $\alpha$, from $0.991$ at $\alpha = 0.1$ to $0.995$ at $\alpha = 5.0$---a spread of $0.004$ in a metric bounded in $[0, 1]$. The high baseline value is geometrically expected (§\ref{sec:measurement-prereqs}). The RepE honesty channel $\rho_R$ varies more substantially, from $-1.306$ at $\alpha = 0.1$ to $-0.533$ at $\alpha = 5.0$, confirming that the linear monitor channel registers the in-direction component of the steered cache state. Cross-channel divergence TDS \emph{decreases} monotonically with $\alpha$ (from $0.094$ to $0.046$), indicating that the H-Neuron and RepE channels agree more strongly under stronger steering---the opposite of the signature predicted by the Orthogonal Escape Hypothesis (Hypothesis~\ref{hyp:orthogonal_escape}). The H-Neuron stress channel $\sigma_H$ reports zero across all $\alpha$; controlled verification (§\ref{sec:measurement-prereqs}) localizes this to a measurement-pipeline issue in the dense-scoring path independent of the steering projection or calibration regime.

\begin{table}[h]
\centering
\small
\begin{tabular}{@{}lrrrrr@{}}
\toprule
$\alpha$ & 0.1 & 0.5 & 1.0 & 2.0 & 5.0 \\
\midrule
OEI & 0.991 & 0.991 & 0.992 & 0.992 & 0.995 \\
$\rho_R$ & $-1.306$ & $-1.258$ & $-1.198$ & $-1.080$ & $-0.533$ \\
$\sigma_H$ & 0.000 & 0.000 & 0.000 & 0.000 & 0.000 \\
TDS & 0.094 & 0.090 & 0.084 & 0.075 & 0.046 \\
\bottomrule
\end{tabular}
\caption{Per-$\alpha$ mean telemetry under inference-time KV-cache steering on Llama~3.2-1B-Instruct, post-projection fix. OEI shows small but monotonic variation in the high-dimensional saturation regime; $\rho_R$ varies substantially, confirming the steering reaches the cache; TDS \emph{decreases}, indicating the channels agree more under stronger steering---the opposite of the Orthogonal Escape signature. $\sigma_H = 0$ throughout reflects monitor calibration prerequisites discussed in §\ref{sec:measurement-prereqs}. Reversibility holds end-to-end (residual\_norm $= 0$).}
\label{tab:exp1-results}
\end{table}

\textbf{Interpretation.} The results are consistent with the negative-result regime Bailey et al.~\cite{bailey2024obfuscated} report for gradient-trained input attacks: high-dimensional null-space saturation dominates, and orthogonal escape does not emerge from non-orthogonality-optimizing search at the model scales tested. We interpret this as a measurement under one specific monitor specification (rank-1 RepE honesty probe at the middle layer); whether richer monitors or hallucination-inducing input regimes produce a more discriminative signal is the natural next experimental step, and is not contradicted by these results.
% [IMPL: COMPLETE — run_kv_mcts.py post-fix, alpha-sweep on Llama-3.2-1B-Instruct, output at tmp/runs/2026-05-06_b6_post-fix.json. 3B reversibility confirmed via probe_kv_cache_mutability.py 2026-05-05.]

% [REVISION | Exp2-recast-reward-ablation | 2026-05-05 | Recast Experiment 2 from "latent signal quality per compute unit vs random KV search" (Option A, 2026-04-25) to a reward-function ablation. Pre-mortem FM 3: Option A inherited the same tautology Option B was rejected for — MCTS optimizes for the very telemetry triple Exp 2 measured, so a random-KV baseline loses by construction. The honest ablation isolates which reward components carry signal: full reward vs single-channel σ_H-only reward vs random-KV baseline. Text-space methods (GCG/MPA/DAMON) move from primary baselines to qualitative discussion in §2. — pending audit]
\paragraph{Experiment 2: Reward-Function Ablation in Latent Space.}
We measure the contribution of each component of the MCTS reward (Eq.~\ref{eq:mcts_reward}) by comparing three search arms operating in the same KV-cache latent space at a fixed compute budget on Llama~3.2-1B-Instruct: (i)~MCTS with the full telemetry reward $r = \lambda_1 \bar{\sigma}_H + \lambda_2 (1-\bar{\rho}_R) + \lambda_3 \, \mathrm{Divergence}(\sigma_H, \rho_R)$; (ii)~MCTS with a $\sigma_H$-only reward (ablating $\rho_R$ and the divergence term, i.e.\ $\lambda_2 = \lambda_3 = 0$); (iii)~random KV-cache mutation (no reward signal). For each arm we record the per-node telemetry triple $(\sigma_H, \rho_R, \mathrm{OEI})$ and report mean and variance across the search budget. This is an honest ablation of the reward formulation rather than a comparison against text-space methods: asking whether a structured search beats unstructured search at finding the structured-search reward is true by construction; the substantive question is which reward components carry signal beyond the simplest single-channel monitor. Existing text-space jailbreak families---GCG~\cite{zou2023gcg}, MPA~\cite{wu2025mcts}, DAMON~\cite{damon2025}---operate on model outputs (text tokens, and in MPA's case first-position target-token log-probabilities) rather than internal representations and therefore produce no per-node latent signal under our metric; we discuss them qualitatively in §2 as the categorical alternative whose per-step compute is invested in text-output evaluation rather than internal-state probing. Cross-method attack-success-rate comparison requires multi-step text rollout and is deferred to Phase~B.
% [IMPL: PARTIAL — run_kv_mcts.py implements the full-reward arm; the σ_H-only-reward arm is a $\lambda$-toggle in the same script (no code change required); the random-KV arm uses the existing random-mutation harness. 2026-05-05 B6 diagnostic surfaced a measurement-pipeline bug producing bit-identical telemetry across $\alpha$; the ablation cannot run until that is resolved (see action plan §4.1.5).]

% [REVISION | §5-Exp3-5-cut | 2026-04-24 | Experiments 3, 4, 5 removed from §5; deferred to Future Work subsection in §7 — pending audit]

\subsection{Baselines}

% ╔══════════════════════════════════════════════════════════════════════════╗
% ║  EDITOR NOTE — TABLE 2: DO NOT ADD PUBLISHED ASR FIGURES FROM BASELINES ║
% ║                                                                          ║
% ║  What this table currently is:                                           ║
% ║    A METHODS REFERENCE TABLE. It lists baselines by category and paper.  ║
% ║    It contains NO experimental results and is NOT a comparison table.    ║
% ║                                                                          ║
% ║  What this table must eventually become (after Phase B experiments):     ║
% ║    A RESULTS TABLE built from measurements we control — all baselines    ║
% ║    run on the same model, same task, same success criterion, same        ║
% ║    payload set. That data comes from Experiment 2 (Phase B, H100).       ║
% ║                                                                          ║
% ║  What NOT to do:                                                         ║
% ║    Do NOT add a published ASR column using figures from the baseline     ║
% ║    papers themselves. GCG's 88% is on Vicuna-7B / AdvBench (discrete    ║
% ║    token jailbreaking). MPA and DAMON figures are from different models  ║
% ║    and benchmarks. Our ASR = "fraction of KV-cache MCTS rollouts that   ║
% ║    produce a successful alignment-breaking intervention" — a different   ║
% ║    task, different model, different success criterion. These numbers      ║
% ║    cannot share a column without becoming a category error.              ║
% ║                                                                          ║
% ║  Before revising this table, verify:                                     ║
% ║    1. Have Phase B experiments completed? (H100, Experiment 2)           ║
% ║    2. Were all baselines run on the SAME model, task, and payload set?   ║
% ║    3. Is the success criterion identical across all rows?                 ║
% ║    If any answer is NO, the table remains a reference table only.        ║
% ║                                                                          ║
% ║  Decision recorded: 2026-04-23 (Josh + Claude)                          ║
% ║                                                                          ║
% ║  UPDATE (2026-04-24): Experiment 2 reframed to Option 2 (latent signal  ║
% ║  quality per compute unit). Eventual results column will compare per-    ║
% ║  node telemetry richness (OEI, σ_H, ρ_R per node evaluation) rather     ║
% ║  than full-response ASR. Update verification checklist above when        ║
% ║  Phase B experiments are complete.                                        ║
% ╚══════════════════════════════════════════════════════════════════════════╝

\begin{table}[h]
\centering
\begin{tabular}{@{}lll@{}}
\toprule
\textbf{Category} & \textbf{Baseline} & \textbf{Reference} \\
\midrule
Behavioral Red-Teaming & GCG (Greedy Coordinate Gradient) & Zou et al., 2023 \\
Text-Space MCTS & MPA (MCTS Prompt Autogeneration) & Wu et al., 2025 \\
Dialogue MCTS & DAMON & 2025 \\
LLM-as-Judge & MT-Bench / G-Eval & Zheng et al., 2023 / Liu et al., 2023 \\
Static Analysis & SonarQube / Semgrep & --- \\
Activation Steering & CAA (Contrastive Activation Addition) & Rimsky et al., 2024~\cite{rimsky2024steering} \\
Cache Steering & One-shot static KV-Cache Steering (precomputed mean-of-differences vector, single post-prefill injection) & Belitsky et al., 2025 \\
Latent Monitor Evasion & RL-Obfuscation & Gupta \& Jenner, 2025 \\
\bottomrule
\end{tabular}
\caption{Baseline methods for comparative evaluation. The Cache Steering row uses Belitsky et al.'s static one-shot paradigm (precomputed mean-of-differences, single injection post-prefill); our MCTS uses iterative reversible search with telemetry-derived rewards---a distinct intervention regime.}
\label{tab:baselines}
\end{table}

\subsection{Evaluation Metrics}

\begin{itemize}[leftmargin=*]
    \item \textbf{Orthogonal Escape Index (OEI):}
    \begin{equation}
    \text{OEI} = \frac{\|\text{Proj}_{\text{Null}(\mathbf{P})}(\mathbf{h}_{\text{steered}} - \mathbf{h}_{\text{base}})\|_2}{\|\mathbf{h}_{\text{steered}} - \mathbf{h}_{\text{base}}\|_2}
    \label{eq:oei}
    \end{equation}
    Fraction of the steering-induced change $\Delta\mathbf{h} = \mathbf{h}_{\text{steered}} - \mathbf{h}_{\text{base}}$ that lies in the null space of the linear monitor $\mathbf{P}$. $\text{OEI} \in [0, 1]$: $\text{OEI} \to 1$ indicates the activation change is entirely orthogonal to the monitor direction (Orthogonal Escape); $\text{OEI} \to 0$ indicates the change is aligned with the monitored direction (normal response).

    % [CRITIQUE NOTE | Metric robustness]
    % OEI may produce false negatives if redistribution is nonlinear or
    % manifold-structured rather than concentrated along linear null-space directions.
    % A model that redistributes deceptive encodings onto a nonlinear manifold within
    % the null space would evade OEI detection while still escaping the linear monitor.
    % Consider acknowledging this in Limitations, and note it as a motivation for
    % future work on nonlinear monitor extensions.

    % [REVISION | T-C-TDS-reframe | 2026-05-05 | Reframed TDS from a standalone evaluation metric to the divergence component of the MCTS reward function. Note the circularity: the same JSD value appears as the third term of Eq.~\ref{eq:mcts_reward} (which the search optimizes) and as the diagnostic signature for the Orthogonal Escape state in Table~\ref{tab:diagnostic_states}; reporting "high TDS at terminal nodes therefore indicates OE" would be tautological. Hence: defined here for completeness, not separately tabulated as a results column in submission experiments. — pending audit]
    \item \textbf{Telemetry Divergence Score (TDS).} The Jensen--Shannon divergence between the H-Neuron stress and RepE projection signal distributions across layers; this is the divergence component referenced in the MCTS reward function (Eq.~\ref{eq:mcts_reward}, $\lambda_3$ term) and is defined here for completeness. Because the search optimizes this quantity directly, TDS is not separately tabulated as an evaluation outcome in submission experiments---reporting elevated TDS at terminal nodes would amount to confirming that an optimizer that maximizes a quantity tends to produce nodes with high values of it. The standalone metric we report from the latent-search process is OEI; TDS surfaces in the rollout traces and per-node telemetry but not as a results-table column.

    % [REVISION | §5.4-metrics-cull | 2026-05-05 | Removed orphaned metric items (ASR, pass@1, Cosine Similarity of Semantic Intent, Memory Efficiency Ratio) and the redundant CRITIQUE NOTE that flagged them. Rationale: Experiments 3-5 were cut on 2026-04-24 (deferred to Future Work); Experiment 2 was reframed to latent signal quality on 2026-04-25 (Option A); none of the four removed metrics are measured by any remaining submission experiment. Surviving inventory in §5.4: OEI (full evaluation metric) + TDS (reward component, defined for completeness). — pending audit]
\end{itemize}

% ============================================================
% 6. THEORETICAL ANALYSIS
% ============================================================
\section{Theoretical Analysis}

\subsection{Reversibility Guarantee}

% [IMPL: EMPIRICALLY VALIDATED — measure_lipschitz_drift.py, n=200 cycles, FP32 accumulator_inf_norm = 0.00e+00]
\begin{theorem}[Exact Reversibility with FP32 Accumulator]
\label{thm:reversibility}
% [REVISION | thm-label-add | 2026-05-06 | Added \label{thm:reversibility} to enable forward references from §1 abstract, §1 Contributions, and §7 Empirical Measurement Prerequisites subsection (all landed 2026-05-06). Theorem statement and proof unchanged. — pending audit]
Let $\mathbf{K}_0 \in \mathbb{R}^{S \times d}$ be the baseline KV cache in bf16. Let $\{\delta_1, \ldots, \delta_n\}$ be a sequence of steering interventions applied and reversed via the FP32 accumulator. Then after $n$ complete apply-reverse cycles:

\begin{equation}
\|\mathbf{K}_n - \mathbf{K}_0\|_\infty \leq n \cdot \epsilon_{\text{bf16}} \cdot \max_i \|\delta_i\|_\infty
\label{eq:reversibility_naive}
\end{equation}

where $\epsilon_{\text{bf16}} = 2^{-8} \approx 3.9 \times 10^{-3}$ is the bf16 machine epsilon. With FP32 accumulation, this bound tightens to:

\begin{equation}
\|\mathbf{K}_n - \mathbf{K}_0\|_\infty \leq \epsilon_{\text{bf16}} \cdot \|\mathbf{A}_{\text{final}}\|_\infty
\label{eq:reversibility_fp32}
\end{equation}

which is independent of $n$, depending only on the final accumulator state (which is zero after complete reversal).
\end{theorem}

\begin{proof}[Proof sketch]
The FP32 accumulator maintains the exact cumulative delta. The only quantization error occurs in the final cast from FP32 to bf16, which is bounded by $\epsilon_{\text{bf16}}$ times the accumulator magnitude. After complete reversal, $\mathbf{A}_{\text{final}} = \mathbf{0}$, so $\mathbf{K}_n = \mathbf{K}_0$ exactly (up to the single bf16 quantization of zero, which is exact). \qed
\end{proof}

% [REVISION | DRAFT-HOLD-block5-consolidated | 2026-05-05 | GAP-C2-03 + GAP-C2-04 (Memory Complexity Proposition: 40.05GB / 9.7TB headline numbers reflect sparse-accumulator + no-clone assumptions; Phase 2 floor is ~2-3·M_KV from full-shape accumulator + baseline clone) — block content moved into §5 Implementation Notes paragraph (item iv). Headline numbers within the proposition itself remain unchanged here pending Max's simplified Exp 3 VRAM measurements (action plan §4.2 task "Apply Memory Complexity Proposition rewrite"). Code fix tracking continues in action plan §4.2 / Phase 3. — pending audit]
\subsection{Memory Complexity}

\begin{proposition}
\label{prop:memory}
% [REVISION | prop-label-add | 2026-05-06 | Added \label{prop:memory} to enable forward references from §1 Contributions and §7 Empirical Measurement Prerequisites subsection (both landed 2026-05-06). Proposition statement unchanged; headline numbers in the comparison sentence below remain TODO pending Max's simplified Exp 3. — pending audit]
The reversible MCTS algorithm achieves memory complexity $O(M_{\text{KV}} + d \cdot K_{\text{acc}})$ where $d$ is the maximum search depth and $K_{\text{acc}} = L \cdot S' \cdot d_{\text{model}} \cdot 4$ bytes is the per-step FP32 accumulator size for $S'$ mutated positions. This is independent of the branching factor $b$.
\end{proposition}

For comparison, standard parallel MCTS requires $O(b^d \cdot M_{\text{KV}})$ memory. With $b = 3$, $d = 5$, and $M_{\text{KV}} = 40\text{GB}$, standard MCTS requires ${\sim}9.7$TB, while reversible MCTS requires ${\sim}40.05$GB.

% ============================================================
% 7. DISCUSSION
% ============================================================
\section{Discussion}

% [REVISION | §7.1-monitor-design-verbs | 2026-05-05 | Reframed §7.1 from attack-implications language ("arms race", "evasion strategies", "alignment assurance against capable models") to monitor-design language (Gemini #1 threat-model critique). The methodological frame: KV-cache MCTS requires white-box access to a frozen model; it is an interpretability / red-teaming tool, not an inference-time attack vector available to a real adversary. Heading retitled to match the verb pivot. Substantive observations preserved (null-space dimensionality, fixed-monitor blind spots, dynamic-monitor opportunity); attack/adversary framing removed. — pending audit]
\subsection{Implications for Latent-Space Monitor Design}

If alpha-sweep characterization surfaces systematic activation redistribution into the null space of a fixed linear monitor, the implication for monitor design is direct: a single static probe direction does not bound the space a model's representations may occupy under intervention. The dimensionality of continuous representation spaces provides a large reservoir of directions orthogonal to any one chosen monitor axis, and a white-box tool that operates inside the frozen model can map this reservoir directly---surfacing where a fixed monitor would produce false negatives without requiring the model to be deceptive in deployment.

The correct frame for these findings is monitor-design diagnostics rather than attack capability: KV-cache MCTS requires white-box access to a frozen model and is therefore unavailable to a real adversary at inference time. Its value is interpretability and controlled-access red-teaming. The forward-looking design implication is that monitor selection should be evaluated against the directions a frozen model's hidden states actually visit under perturbation, not only against the directions a probe was trained on. The reversible MCTS framework presented here provides the computational infrastructure for such evaluation; the same algorithm, paired with an adaptive-probe reward, would support monitor-selection procedures that update the monitored subspace in response to observed redistribution patterns surfaced by the search.

\subsection{Broader Impact}

This work has dual-use implications that we address transparently. The Reversible MCTS framework could, in principle, be used to discover alignment-breaking interventions that are then exploited rather than patched. We mitigate this risk through three design choices: (1)~the framework operates in a controlled, offline environment rather than against live systems; (2)~a cryptographic hash of the evaluation payload provides an audit trail of all interventions and discoveries; (3)~we commit to responsible disclosure of any critical vulnerabilities discovered during experimentation.

\subsection{Limitations}

We acknowledge several limitations. First, the Orthogonal Escape hypothesis, while motivated by the RL-Obfuscation findings~\cite{gupta2025rlobfuscation}, has not yet been empirically validated in the specific setting of KV-cache interventions; our Phase~A experiments are designed to provide this validation. Second, the cross-model transfer of steering vectors via Procrustes alignment is an open research question with uncertain guarantees. Third, the framework relies on H-Neuron activation tracking, which Gao et al.\ (2025) associated with hallucination and over-compliance; the assumption that these neurons also serve as reliable indicators of complex, instrumentally rational strategic deception remains an open theoretical liability that our experiments must address. Additionally, Zou et al.\ (2023) validated RepE on models of 7B parameters and above; our 1B surrogate falls outside this validated range, and probe reliability at sub-7B scale is an open empirical question. Fourth, the current experimental design would benefit from cross-family scaling validation across architectures beyond the Llama family to establish generalizability of the Orthogonal Escape phenomenon. Fifth, $\sigma_H$ may conflate hallucination pressure, mode-collapse, and strategic masking under corrective pressure unless discriminant controls are applied; the extension from hallucination association to instrumentally rational deception is a testable assumption rather than a settled equivalence. Sixth, OEI may produce false negatives if deceptive redistribution is nonlinear or manifold-structured rather than concentrated along linear null-space directions, in which case $\text{OEI}\to 0$ while the model still evades the linear monitor. Finally, the proposed framework requires an explicit validation chain linking continuous latent interventions to discrete deceptive-intent outcomes; Phase~B experiments are designed to provide this evidence but have not been run at time of submission.
% [REVISION | G1 | 2026-04-24 | Claims 12+13+14 added to §7 Limitations — pending audit]

% [REVISION | §7-measurement-prereqs | 2026-05-06 | Inserted new §7 subsection "Empirical Measurement Prerequisites for OEI" between §7.3 Limitations and §7.4 Future Work (Session B's structural placement choice — see §7 reply in action plan). Lands \label{sec:measurement-prereqs} which the abstract and §1 Contributions piece #2 reference. Honest finding documents the GAP-C2-06 measurement-pipeline issue (per Option α+ resolution in pre-mortem decision 2). Three resolution approaches outlined; Theorem 1 + memory bound explicitly stated as unaffected. — pending audit]
% [REVISION | §7.4-post-fix-update | 2026-05-06 | Replaces the morning's draft of this subsection (§7-measurement-prereqs marker above retained as the original insertion record). Tightened to reflect Josh's verified GAP-C2-06 fix landed 2026-05-06 PM: removed "early version" placeholder language; the W_K projection is now the realized fix, not a candidate. References §\ref{sec:experiments} for results numbers (rather than restating). Added a new third paragraph documenting σ_H = 0 as a calibration-prerequisite finding (per Session A's analysis: H-Neuron monitor is contrastively calibrated against hallucination-inducing prompts, our test prompts are not hallucination-inducing, per-neuron ReLU correctly clamps below-baseline activations). Section heading retitled "Empirical Measurement Prerequisites and Initial Observations" to reflect the addition of measured findings. — pending audit]
\subsection{Empirical Measurement Prerequisites and Initial Observations}
\label{sec:measurement-prereqs}

The Orthogonal Escape Index in Eq.~\ref{eq:oei} requires that the steering direction $\mathbf{d}_K^{(\ell)}$ live in the same space as the cached key tensors. Representation-engineering probes are calibrated against residual-stream activations of dimension $d_{\text{model}}$, while the KV cache stores per-head representations of dimension $\text{kv\_heads} \times d_{\text{head}}$ under grouped-query attention. For Llama~3.2-1B these are 2048 and $8 \times 64$ respectively; an explicit projection between probe space and cache space is required. In our reference implementation we identified, characterized, and resolved a measurement-pipeline failure mode in which the broadcast layer silently zeroed the steering delta on shape incompatibility, producing forward passes that were bit-identical across $\alpha$. The fix projects probe directions through layer-specific key projection weights, $\mathbf{d}_K^{(\ell)} = \mathbf{W}_K^{(\ell)} \mathbf{d}_{\text{probe}}^{(\ell)}$, reshaped to per-head form---interpretable as the cache state that would result from the corresponding residual-stream perturbation modulo position-dependent rotary embedding. Two alternative approaches are also viable: intervening at the residual stream prior to the K/V projection and re-running the layer's forward computation, or calibrating probes directly in per-head cache space against contrastive cache states. Systematic comparison of the three is left to future work.

Under the projection-based fix, OEI varies monotonically with steering magnitude on Llama~3.2-1B-Instruct (results in §\ref{sec:experiments}). The high baseline value (${\sim}0.99$) is geometrically expected: in $d_{\text{model}} = 2048$ with a rank-1 linear monitor, removing one direction from any perturbation leaves ${\sim}99\%$ of the perturbation's norm in the null space by default. This high-dimensional saturation is consistent with the regime Bailey et al.~\cite{bailey2024obfuscated} characterize when discussing why activation cosine similarity is hard to drive low in a behavior-preserving setting; the small but monotonic OEI spread we observe suggests the metric is responsive to genuine steering structure rather than producing a uniform default. Whether OEI under a richer monitor (rank-$k$ probe matrix, or attention-aggregated probes; cf.~\cite{gupta2025rlobfuscation}) yields a more discriminative signal is an open empirical question and a natural extension of the framework.

% [REVISION | §7.4-σH-verification-update | 2026-05-06 | Replaces the σ_H paragraph (third paragraph of §7.4) following the diagnostic session's controlled three-prompt verification (recorded in `docs/logs/2026-05-05_diagnostic-report.md` § "σ_H VERIFICATION FOLLOW-UP — 2026-05-06"). The earlier H_A "calibration-prerequisite" framing is refuted: σ_H = 0 across all 17 layers under all three regimes (hallucination / coercion / neutral), including a coercion-style prompt drawn verbatim from the `_COERCED_EXAMPLES` set used as `hallucinated_examples` during H-Neuron calibration. ρ_R varied substantially (deepest-layer −14.66 / −7.15 / −2.93), confirming hidden-state delivery is intact and localizing the constancy to the H-Neuron dense-scoring path itself. New paragraph identifies this as a *second* measurement-pipeline issue (distinct from and independent of GAP-C2-06 resolved by the W_K projection) and leaves resolution as future work. Theorem 1 + Memory Proposition reassurance is folded into this same paragraph (no separate paragraph 4 needed). Two earlier markers retained as layered audit trail: `§7-measurement-prereqs` (original AM insertion) and `§7.4-post-fix-update` (afternoon update). — pending audit]
% [REVISION | §7.4-σH-investigation-update | 2026-05-07 | Additive update to §7.4 paragraph 3 following the diagnostic session's GAP-C1-02 investigation (`docs/logs/2026-05-07_optional-σH-investigation-prompt.md`). Path A confirmed GAP-C1-02 (per-neuron baseline-subtraction-and-ReLU) is not applied in current code; Path B applied the fix tentatively per Max's recipe and re-ran the three-prompt verification — σ_H remained identically zero across all 17 layers in all three regimes (ρ_R / OEI / TDS bit-identical to yesterday's verification, confirming the tentative fix touched only the σ_H pipeline). Cause therefore localized as: downstream of per-neuron normalization. Most plausible remaining candidate: position/context discrepancy between calibration-time (last-input-token, prompt-only forward pass) and inference-time (last-token, prompt-plus-generation, steered cache) hidden states. Replaces the "Constancy of σ_H..." sentence with two new sentences capturing the GAP-C1-02 ruling-out and the position/context hypothesis; rest of paragraph (opening sentences + future-work conclusion + Theorem 1/Proposition reassurance) preserved verbatim from yesterday's drafts v3 patch. Net +60 words. Three prior §7.4 markers retained as layered audit trail: `§7-measurement-prereqs` (2026-05-06 AM), `§7.4-post-fix-update` (2026-05-06 PM), `§7.4-σH-verification-update` (2026-05-06 EOD). — pending audit]
The bottom-up H-Neuron stress channel ($\sigma_H$) is reported as zero across all $\alpha$ in our reference experiments. This is independent of the steering pipeline (which demonstrably reaches the cache, per the per-$\alpha$ variation in $\rho_R$, TDS, and OEI in §\ref{sec:experiments}). To distinguish whether the constancy reflects a calibration-regime mismatch~\cite{gao2025hneurons} or a deeper measurement issue, we ran a controlled three-prompt comparison at fixed $\alpha = 1.0$, varying only the test prompt across hallucination-style, coercion-style, and neutral regimes---with the coercion-style prompt drawn directly from the same class used as the H-Neuron monitor's contrastive calibration set. $\sigma_H$ remained identically zero across all 17 layers in all three regimes, while $\rho_R$ varied substantially (deepest-layer values of $-14.66$, $-7.15$, and $-2.93$ across the three regimes respectively). To narrow the cause, we additionally tested the per-neuron baseline-subtraction-and-ReLU formulation specified by Eq.~\ref{eq:hneuron_stress} (which our reference implementation aggregates rather than applying per-neuron); under this formulation $\sigma_H$ remained identically zero across the same three regimes, indicating the cause is downstream of per-neuron normalization. The remaining most plausible cause is a position/context discrepancy between calibration-time and inference-time hidden states: the H-Neuron monitor is calibrated on last-input-token activations from prompt-only forward passes, while the MCTS evaluation reads last-token activations from prompt-plus-generation forward passes under a steered cache; these are different positions in the sequence with different attention contexts, and the H-Neurons selected by the former calibration may not fire under the latter regime even when input semantics are similar. We identify this as a second measurement-pipeline issue, and validation of the bottom-up channel as a precondition for full telemetry-matrix evaluation; resolution is left as future work. We note that Theorem~\ref{thm:reversibility} and Proposition~\ref{prop:memory} are unaffected by these measurement choices: both are intervention-agnostic.

\subsection{Future Work}
% [CRITIQUE NOTE | Future Work section — NeurIPS convention check]
% Verify whether it is standard practice for NeurIPS Main Track papers to include a dedicated
% Future Work subsection describing planned (unrun) experiments. Some venues prefer this content
% folded into Limitations or Discussion rather than as a standalone subsection.
% Confirm with Tianyu Shi or check accepted NeurIPS 2024/2025 papers in the alignment/interpretability area.

Several directions are deferred to Phase~B. Experiments~3--5 (memory efficiency profiling, evaluation reproducibility, and cross-model steering transfer) require infrastructure not yet complete: sparse FP32 accumulators, a ground-truth execution sandbox, and Procrustes alignment between model families. Scaling to \texttt{openai/gpt-oss-20b} (21B parameters, Apache~2.0) will test whether the Orthogonal Escape phenomenon and telemetry diagnostics hold at 20B scale, with router logit entropy as the H-Neuron proxy for MoE layers. Full sparse accumulator implementation---reducing per-step overhead from $O(S \cdot d_{\text{model}})$ to $O(S' \cdot d_{\text{model}})$ for $S' \ll S$ mutated positions---is required before the memory complexity proposition can be validated empirically. A Croissant~1.1 dataset package of runtime evidence artifacts is planned for release alongside the Phase~B submission.
% [REVISION | §7-FutureWork | 2026-04-24 | Future Work subsection added covering Experiments 3–5, gpt-oss-20b Phase B, sparse accumulators — pending audit]
% [REVISION | §7-Croissant | 2026-04-25 | Appended Croissant 1.1 runtime-evidence package sentence (ORIGINAL PLAN omission); complementary to §7 Ethical Considerations Croissant metadata sentence — pending audit]

\subsection{Ethical Considerations}

All experiments will be conducted on open-weight models in isolated sandbox environments. No experiments will target production systems or real users. The adversarial datasets generated will be released under a responsible disclosure framework with appropriate access controls. The release package will include a Croissant 1.1 metadata file (plus Croissant RAI extension fields) so provenance, usage constraints, and schema-level reproducibility checks are machine-verifiable.

% [REVISION | §8-cut | 2026-04-24 | Timeline and Milestones section removed entirely — research proposal artifact, not a paper section; ~0.5 page saved — pending audit]

% ============================================================
% REFERENCES
% ============================================================
\begin{thebibliography}{22}

\bibitem{zheng2023judging}
L.~Zheng, W.-L.~Chiang, Y.~Sheng, et al.
\newblock Judging LLM-as-a-Judge with MT-Bench and Chatbot Arena.
\newblock In \emph{NeurIPS}, 2023.

\bibitem{liu2023geval}
Y.~Liu, D.~Iter, Y.~Xu, et al.
\newblock G-Eval: NLG Evaluation using GPT-4 with Better Human Alignment.
\newblock In \emph{EMNLP}, 2023.

\bibitem{greenblatt2024alignment}
R.~Greenblatt, C.~Shlegeris, et al.
\newblock Alignment Faking in Large Language Models.
\newblock Anthropic Research, December 2024.

\bibitem{hubinger2024sleeper}
E.~Hubinger, C.~Denison, J.~Mu, et al.
\newblock Sleeper Agents: Training Deceptive LLMs That Persist Through Safety Training.
\newblock \emph{arXiv:2401.05566}, 2024.

\bibitem{gupta2025rlobfuscation}
R.~Gupta and E.~Jenner.
\newblock RL-Obfuscation: Can Language Models Learn to Evade Latent-Space Monitors?
\newblock \emph{arXiv:2506.14261v4}, 2025.

\bibitem{gao2025hneurons}
C.~Gao et al.
\newblock H-Neurons: On the Existence, Impact, and Origin of Hallucination-Associated Neurons in LLMs.
\newblock \emph{arXiv:2512.01797}, 2025.

\bibitem{zou2023repe}
A.~Zou, L.~Phan, S.~Chen, et al.
\newblock Representation Engineering: A Top-Down Approach to AI Transparency.
\newblock \emph{arXiv:2310.01405}, 2023.

\bibitem{zou2023gcg}
A.~Zou, Z.~Wang, J.~Z.~Kolter, and M.~Fredrikson.
\newblock Universal and Transferable Adversarial Attacks on Aligned Language Models.
\newblock \emph{arXiv:2307.15043}, 2023.

\bibitem{wu2025mcts}
S.~Wu et al.
\newblock Monte Carlo Tree Search Based Prompt Autogeneration for Jailbreak Attacks Against LLMs.
\newblock In \emph{COLING}, 2025.

\bibitem{damon2025}
DAMON: A Dialogue-Aware MCTS Framework for Jailbreaking Large Language Models.
\newblock In \emph{EMNLP}, 2025.

\bibitem{li2025ace}
X.~Li et al.
\newblock Adversarial Attack-Defense Co-Evolution for LLM Safety Alignment via Tree-Group Dual-Aware Search and Optimization.
\newblock \emph{arXiv:2511.19218}, 2025.


\bibitem{turner2023steering}
A.~M.~Turner, L.~Thiergart, G.~Leech, D.~Udell, et al.
\newblock Steering Language Models with Activation Engineering.
\newblock \emph{arXiv:2308.10248}, 2023.

\bibitem{belitsky2025kvcache}
M.~Belitsky, D.~J.~Kopiczko, M.~Dorkenwald, et al.
\newblock KV Cache Steering for Controlling Frozen LLMs.
\newblock \emph{arXiv:2507.08799}, 2025.

\bibitem{agenticred2025}
AgenticRed: Optimizing Agentic Systems for Automated Red-teaming.
\newblock \emph{arXiv:2601.13518}, 2025.

\bibitem{llmasjudge2025limitations}
Limitations of LLM-as-a-Judge Without Human Grounding.
\newblock \emph{arXiv:2503.05061}, 2025.


\bibitem{hao2025coconut}
S.~Hao, B.~Suber, Z.~Wang, et al.
\newblock Training Large Language Models to Reason in a Continuous Latent Space.
\newblock In \emph{ICLR}, 2025. \emph{arXiv:2412.06769}.

% [REVISION | G3 | 2026-04-24 | crsm2025 bibitem removed entirely — unreviewed prototype, no peer-reviewed venue — pending audit]

\bibitem{bailey2024obfuscated}
J.~Bailey et al.
\newblock Obfuscated Activations Bypass LLM Latent-Space Defenses.
\newblock In \emph{NeurIPS}, 2024.

\bibitem{rimsky2024steering}
N.~Rimsky, N.~Gabrieli, J.~Schulz, M.~Tong, E.~Hubinger, and A.~Turner.
\newblock Steering {L}lama 2 via Contrastive Activation Addition.
\newblock In \emph{Proceedings of the 62nd Annual Meeting of the Association for Computational Linguistics (Volume 1: Long Papers)}, pages 15504--15522, Bangkok, Thailand, 2024. Association for Computational Linguistics.
% [REVISION | G2 | 2026-04-24 | panickssery2023steering replaced with rimsky2024steering (ACL 2024) — pending audit]

\bibitem{okazakirag2025}
System 2 through Okazaki-RAG: Capacity-Bounded Inline Fragment Manipulation for Deliberative Reasoning.
\newblock ResearchGate, 2025.

\end{thebibliography}

\end{document}
